\chapter[The Conclusion]{The Conclusion}
\label{Chap7}
%=============================================================================================================
\lhead{\emph{Chapter 7: The Conclusion}} % Write in your own chapter title to set the page header
\epigraph{\itshape ``Statistics: The only science that enables different experts using the same figures to draw different conclusions.''}{-- Evan Esar
}



\noindent

Now, I came to the end of the thesis, which brings me to discuss the conclusion, novelty, and future aspects of it. In the entire thesis, I have primarily used the white light data from KoSO along with white light and magnetogram data from space-based modern instruments MDI and HMI. With the help of these data sets, I have explored the multiple observational aspects of the solar long--term variability and how it can reduce the gap between observation and the theoretical aspects of solar dynamo models. Let me start to conclude my thesis for each chapter separately.

\section{Chapter~3}

Motivated by the work of \citet[references therein]{Hathaway2013} and understanding the importance of penumbra to umbra area ratio ($q$), I have studied the long-term variation of this ratio using the historical KoSO digitized data (1923--2011). I have developed an automatic method based on Otsu technique to extract the penumbra and umbra from the sunspots. I have noted that the $q$ increases initially for area $<200\mu$Hem, but after that, it increases very slow and saturates around 5.5 to 6, which is in agreement with earlier findings. Although the bigger sunspots (Area$>200\mu$Hem) show a similar trend in both the data sets (RGO \& KoSO), the smaller sunspots (sunspot area $<100~\mu$Hem) in KoSO data does not show any systematic variation in the ratio contrary to the result reported by \citet{Hathaway2013}. It has been further verified using MDI white--light data \citep{Jha2019}. The result from RGO data raise a question about the origin of two classes of sunspots, but since KoSO data does not show any difference in both, the classes imply they are from the same source. Hence, dynamo modellers might treat them on the same footing.

\section{Chapter~4}

The significant role of solar differential rotation and its variation for the solar dynamo models; and the unavailability of an automated tracking algorithm to track the sunspots motivated me to measure it using a computerized method. I have developed an automatic way to track the sunspots from the century-long KoSO digitized white-light data and calculated the solar rotation profile. Average rotation profile measured in \cref{Chap4} show a very good agreement with the earlier results \citep{Gupta1999} with the following rotation profile when measured in degree/day,
%
\begin{equation}
    \Omega (\theta) = (14.381\pm 0.004)-(2.72\pm0.04)\cos^2{\theta},
\end{equation}
where $\theta$ is the co-latitude. KoSO data neither show any phase dependency nor the significant long--term variation in the solar differential rotation. However, differential rotation parameters $A$ and $B$ show the negative and positive correlation with the strength of the solar cycle. In the case of KoSO white--light data, the smaller sunspots (area $<200~\mu$Hem) gives a faster rotation rate than the bigger ones (area $>400~\mu$Hem) which is believed to be due to their different anchored depth and drag imposed by them. 

\section{Chapter~5}
%
The differential rotation measurement based on the Doppler velocity \cite{Howard1970} gives a slightly lower rotation rate when compared with the one measured using sunspots \citep[\cref{Chap4}, and ][]{Beck2000, Jha2021}.  This is due to the sharp change in the solar differential profile near the surface called NSSL, and the helioseismology observation has confirmed the existence of this intriguing layer. The absence of theoretical understanding of this layer made me think about its physics. Therefore, based on the thermal wind balance equation, which tells how the slight difference in temperature between solar pole and equator is balanced by centrifugal force, a theoretical model for NSSL is presented. Here the idea is the Coriolis force, arising due to solar rotation, has an effect on the convective blobs in the inner convection zone where convective turnover time ($\tau$) is comparable (or greater) than the solar rotation period ($T$) whereas it will be unaffected, near the surface ($\tau<<T$ ). Hence, near the surface the rate of fall of temperature becomes independent of latitude making the thermal wind term very large. The centrifugal term increases to balance the increase in the thermal wind term, leading to the observed NSSL. The analytical expression of $\Omega (r,\theta) $ and the observed $\Omega (r, \theta)$ gives $r_c=0.92$~\Rsun\ and 0.96~\Rsun\ respectively. Here convection zone model, which is called MODEL-S \citep{Dalsgaard1996} is used for the temperature profile. The result based on the observed profile of solar differential rotation and corresponding calculated PETD show good agreement with measured PETD in \citet{Rast2008}.



\section{Chapter~6}

Tilt quenching has been regularly used in the different types of solar dynamo models, particularly the \bl\ type dynamos. What these models were lacking was the observation support to this idea. Although the \citet{DasiEspuig2010} has used the white light data, which do not have the magnetic filed information, and given some indication of the tilt quenching but the measurement based on the magnetogram data was not available. Therefore, in \cref{Chap6} I have first developed and automatic method to identify the BMRs from the MDI (1996--2011) and HMI (2010--2018) magnetogram, which is inspired from the method described in \citet{Stenflo2012a}. Interestingly, in this work a bimodal distribution of \Bmax\ has been found where each peaks corresponds to the two classes of BMR; one which does not show the sunspot (limited by resolution and contrast) in white light (peaks around 600~G) and second which show sunspot in white light (peaks around 2100~G). The \Bmax\ dependence of amplitude of Joy's law ($\gamma_0$) shows an initial increase in $\gamma_0$ and then after around 2~kG it start decreasing and this behaviour is present in both MDI and HMI data sets. This result gives an indication of tilt quenching which need to be further verified using longer data sets with stronger solar cycle.


\section{Novelty of the Thesis}

$\bm{\Rightarrow}$ The first time, a fully automatic method is developed to extract the penumbra, and umbra area from KoSO digitized white--light data. The results obtained in this work does not show any difference in the long--term variation in penumbra to umbra area ratio as speculated based on RGO data (\cref{Chap3}).

$\bm{\Rightarrow}$ A new and automated algorithm to track the sunspots has been developed and used to track the sunspot in KoSO white--light digitized data to measure the solar differential rotation. A significant difference in measurement of solar differential rotation when bigger and smaller sunspots are used has been reported (\cref{Chap4}).

$\bm{\Rightarrow}$ A theoretical model for the NSSL based on thermal wind balance equation, which shows an excellent agreement with observation,  has been given very first time (\cref{Chap5}).

$\bm{\Rightarrow}$ The indication of tilt quenching based on the magnetogram data from MDI and HMI along with the two classes of BMRs has been found, which was not reported earlier. 

\section{Future prospects}

This thesis primarily focused on developing several new automated algorithms to extract different magnetic features as seen on the solar surface, study them, and provide input to the solar dynamo models. There is always a possibility of improving these techniques, and the recent development in computational power, Machine Learning (ML) and Artificial Intelligence (AI) could make it even more efficient and sophisticated. Here I will discuss a few possibilities that I am currently exploring and looking for.

\begin{enumerate}
    \item In \cref{Chap3}, I have used the data from KoSO and MDI, which can be further extended to data from other observatories such as the National Astronomical Observatory of Japan (NAOJ) and others. It will help us to test the algorithm with different data set to improve our statistics and consistency.
    
    \item The method developed in \cref{Chap4} only tracks the sunspots in two consecutive observations, and hence it can be modified to follow the sunspots for full-disk passage or their lifetime. It will help us get more accurate results for the solar differential rotation.
    
    \item A theoretical model of NSSL, given in \cref{Chap5}, also predict the variation of solar surface temperature from the pole to the equator, which can be verified by future observation with the help of more accurate observations.
    
    \item To understand the magnetic field dependence of tilt quenching in \cref{Chap6} I have counted each BMR multiple times as every magnetogram has been treated independently. Therefore,  to better understand the dependence of tilt on \Bmax\ in BMRs it is crucial to track all of them for disk passage or their lifetimes. I am currently involved in the work that is exploring this idea.
\end{enumerate}
